\documentclass[a4paper,10pt]{article}
\newcommand\subdir{/home/koen/LaTeX-setup/}
\input{\subdir/LaTeX-files/imports.tex}
\input{\subdir/LaTeX-files/master-internship.tex}
\usepackage{\subdir/LaTeX-files/aastex}
\usepackage{natbib}
\bibliographystyle{\subdir/LaTeX-files/astroadstitle.bst}

%Set the title, Author and date
\title{Notes Week 37}
\author{Koen Essers}
\tutorial{Master Internship}
\date{\today}

%Set the header style
\header{\tutorialstring}

%Begin the document
\begin{document}
\setuplayout
\section{Increasing the mass-radius grid}

\begin{figure}[ht]
	\centering
	\incplot{w37-roche-radii-2}
\end{figure}

\newsection{Mass-Radius grid}

Below I show some results of the \(M\)--\(R\) grid.
This grid used \(\varepsilon =0.1\), \(\delta = 0.1\),
\(\beta =0.8\) and \(q=0.9\).

To start, let me plot the final periods of this grid.
\begin{figure}[ht]
	\centering
	\incplot{w37-mass-3}
\end{figure}

\begin{figure}[ht]
	\centering
	\incplot{w37-mass-4}
\end{figure}

\begin{figure}[ht]
	\centering
	\incplot{w37-mass-5}
\end{figure}

\begin{figure}[ht]
	\centering
	\incplot{w37-mass-6}
\end{figure}

The RLOF region is clearly behaving very different from the wind region.
We see that the RLOF region has the largest increase in period for small initial periods and large initial masses.

The increase with mass is quite easy to understand, as \(\delta =0.1\) and \(\beta =0.8\) will generally increase the \(P_\text{f} / P_\text{i}\), and a larger mass means the final
\(q\) will be larger.
The change with initial period is an effect of the wind.
I can actually reproduce this using just the analytic formulations of \citet{soberman1997} and the fraction of mass that is lost through
winds.

Below I plot this fraction.

\begin{figure}[ht]
	\centering
	\incplot{w37-mass-8}
\end{figure}

If we take this into account and give the proper fraction of \(\alpha , \beta \) and \(\delta \) to the \citet{soberman1997} formulation,
we obtain the following plot:

\begin{figure}[ht]
	\centering
	\incplot{w37-mass-9}
\end{figure}

The values are not identical, but this is expected since the \citet{soberman1997} prescription is a very simple one.
One source of errors that I can think of are the tides, which are not taken into account
in the \citet{soberman1997} prescription but generally \emph{decrease} the periods before interaction.
This could explain why the analytic predictions are larger.

\nextpage
The wind region in the \(P_\text{f} / P_\text{i}\) plot is actually interesting too.
We see the transition to isotropic winds happening for larger
initial periods.

The figure below shows this the isotropic wind regime better.
\begin{figure}[ht]
	\centering
	\incplot{w37-mass-7}
\end{figure}

Clearly, for much larger periods, the period dependence drops out and we are left with a pure mass dependence.
This mass dependence is caused by the changing final mass ratio.
All models start with \(q=0.9\), but
stars with higher masses can lose more envelope and thus reach larger final \(q\) values, which increases the
fraction \(P_\text{f} / P_\text{i}\).

\begin{figure}[ht]
	\centering
	\incplot{w37-mass-10}
\end{figure}

\bibliography{master.bib}
\end{document}
